\documentclass[a4paper,12pt]{article}
\usepackage[T2A]{fontenc}
\usepackage[utf8]{inputenc}
\usepackage[russian]{babel}
\usepackage{amsmath,amssymb,mathtools}
\usepackage{geometry}
\geometry{left=20mm,right=20mm,top=20mm,bottom=23mm}
\usepackage{helvet}
\renewcommand{\familydefault}{\sfdefault}
\usepackage{microtype}
\usepackage{array,booktabs,tabularx}
\usepackage{tikz}
\usepackage{xcolor}
\usepackage{fancyhdr}
\usepackage[unicode=true,colorlinks=true,linkcolor=darkaccent]{hyperref}
\definecolor{accent}{HTML}{087C83}
\definecolor{darkaccent}{HTML}{07595D}
\definecolor{ink}{HTML}{27333B}
\definecolor{wrong}{HTML}{A53639}
\color{ink}
\setlength{\parindent}{0pt}
\setlength{\parskip}{5pt}
\linespread{1.13}
\pagestyle{fancy}
\fancyhf{}
\renewcommand{\headrulewidth}{0.25pt}
\fancyhead[L]{\scriptsize \textcolor{darkaccent}{МАТЕМАТИКА / РЕВЬЮ ДОМАШНЕЙ РАБОТЫ}}
\fancyhead[R]{\scriptsize 10.11.2026}
\fancyfoot[C]{\scriptsize Лёвин Артём Александрович эксклюзивно для Николь Саркисьянц}\fancyfoot[R]{\scriptsize\thepage}
\newcommand{\errorlabel}{\textcolor{wrong}{\textbf{Ошибка.}}}
\newcommand{\keyidea}{\textcolor{darkaccent}{\textbf{Ключевая идея.}}}
\newcommand{\result}{\textcolor{darkaccent}{\textbf{Итог.}}}
\newcommand{\subline}[1]{\par\vspace{3pt}{\color{darkaccent}\bfseries #1}\par}
\begin{document}
\thispagestyle{empty}
\begin{center}
\vspace*{28mm}
{\color{darkaccent}\large МАТЕМАТИКА}\par\vspace{15mm}
{\fontsize{23}{27}\selectfont Ревью домашней работы}\par\vspace{11mm}
{\fontsize{17}{23}\selectfont\color{accent}Логарифмы и максимум мощности}\par\vspace{15mm}
{\small Подготовка к ЕГЭ по профильной математике}\par
\vfill
{\normalsize Николь Саркисьянц}\par\vspace{8mm}
{\normalsize 10 ноября 2026 г.}\par\vspace{12mm}
{\small Лёвин Артём Александрович}\par
{\small эксклюзивно для Николь Саркисьянц}\par\vspace{12mm}
\begin{tikzpicture}\draw[accent,line width=.9pt] (0,0) .. controls (1.1,.35) and (2.1,-.35) .. (3.2,0) .. controls (4.4,.35) and (5.7,-.35) .. (7.0,0);\end{tikzpicture}
\end{center}
\clearpage
\section*{Сводная таблица}
\textcolor{darkaccent}{Источник:} «Ревью вариант 09.11.26.pdf», упражнения к заданиям 16 и 17.

В таблицу включены только задания, в которых по повторной проверке обнаружены ошибки. Номер задания ведёт к подробному разбору.
\par\par\vspace{4mm}
\small\renewcommand{\arraystretch}{1.35}
\begin{tabularx}{\linewidth}{@{}p{12mm}p{27mm}X p{32mm}p{35mm}@{}}
\toprule
\textbf{№} & \textbf{Тема} & \textbf{Суть ошибки} & \textbf{Ваш ответ} & \textbf{Верный ответ}\\\midrule
\hyperref[task:16]{16} & Логарифмы & В ответ включена точка $x=2$, запрещённая ОДЗ. & \shortstack[l]{$(-\infty;-6]$\\$\cup[2;+\infty)$} & \shortstack[l]{$(-\infty;-6]$\\$\cup(2;+\infty)$}\\
\hyperref[task:17]{17} & Производная & Пропущен квадрат при переходе от $R^2=4$ к $R=4$; итог впоследствии исправлен. & $R=2$, $P_{\max}=288$ & $R=2$, $P_{\max}=288$\\
\bottomrule
\end{tabularx}
\normalsize
\clearpage
\thispagestyle{empty}
\vspace*{\fill}
\begin{center}{\fontsize{25}{30}\selectfont\color{darkaccent}Разбор ошибок}\par\vspace{5mm}
{\large Подробное объяснение и одна задача на закрепление}\end{center}
\vspace*{\fill}
\clearpage
\phantomsection\label{task:16}
\section*{Задание № 16}

\subline{Условие}
Решите неравенство
\[
\log_{0{,}5}\!\left(\frac{x-2}{x+2}\right)+2\log_4(x^2-4)\ge4.
\]
\subline{Ваше решение}
Вы правильно перешли к основанию $2$, объединили логарифмы и получили
\[
\log_2(x+2)^2\ge4 \quad\Longleftrightarrow\quad (x+2)^2\ge16.
\]
Затем решили квадратное неравенство и записали два луча, включив обе граничные точки.

\textbf{Ваш ответ:} $(-\infty;-6]\cup[2;+\infty)$.

\subline{Где именно возникает ошибка}
\textbf{Последний правильный шаг:} $(x+2)^2\ge16$, откуда $x\le-6$ или $x\ge2$.

\errorlabel\ В окончательный ответ включено значение $x=2$. Оно запрещено исходной областью допустимых значений, поскольку при $x=2$ оба исходных аргумента логарифмов обращаются в ноль.

\subline{Почему этот шаг неверен}
Преобразования логарифмов допустимы только на множестве, где все исходные логарифмы определены. При сокращении множителя $x-2$ в дроби мы временно упростили выражение, но запрет на $x=2$ сохранился. При замене выражения более простым нельзя расширять исходную область определения.

\keyidea\ Найдите ОДЗ перед преобразованиями и пересеките её с решением полученного простого неравенства. Закрашенная точка в ответе допустима только при выполнении исходных ограничений.

\subline{Исправление с последнего правильного шага}
\[
(x+2)^2\ge16 \iff x\le-6\quad\text{или}\quad x\ge2.
\]
Из исходных логарифмов получаем $x<-2$ или $x>2$. Следовательно,
\[
\boxed{x\in(-\infty;-6]\cup(2;+\infty)}.
\]
\subline{Полное правильное решение}
\textbf{1. Установим область допустимых значений.}
\[
\frac{x-2}{x+2}>0,\qquad x^2-4>0.
\]
Первое условие выполняется при $x<-2$ или $x>2$; второе --- при тех же ограничениях. Итак,
\[
\text{ОДЗ}:\quad x\in(-\infty;-2)\cup(2;+\infty).
\]
\textbf{2. Приведём логарифмы к основанию $2$.}
\[
\log_{0{,}5}A=-\log_2A,\qquad 2\log_4B=\log_2B\quad (A,B>0).
\]
Поэтому на ОДЗ:
\begin{align*}
-\log_2\!\left(\frac{x-2}{x+2}\right)+\log_2(x^2-4)&\ge4,\\
\log_2\!\left(\frac{(x^2-4)(x+2)}{x-2}\right)&\ge4,\\
\log_2(x+2)^2&\ge4.
\end{align*}
Сокращение правомерно, поскольку на ОДЗ $x\ne2$ и $x\ne-2$.

\textbf{3. Решим полученное неравенство.} Так как основание $2$ больше единицы,
\begin{align*}
(x+2)^2&\ge16,\\
|x+2|&\ge4,\\
x+2&\le-4\quad\text{или}\quad x+2\ge4,\\
x&\le-6\quad\text{или}\quad x\ge2.
\end{align*}
\textbf{4. Пересечём с ОДЗ.} Левый луч целиком допустим. Из правого луча удаляем точку $x=2$.

\subline{Проверка}
При $x=-6$ исходные аргументы положительны, $(x+2)^2=16$, поэтому достигается равенство. При $x=3$ исходные логарифмы определены, $(3+2)^2=25>16$, точка подходит. При $x=2$ обе логарифмические функции в условии не определены; поэтому граница справа открыта.

\result\ $\displaystyle x\in(-\infty;-6]\cup(2;+\infty)$.

\subline{Задача на закрепление}
Решите неравенство
\[
\log_{0{,}5}\!\left(\frac{x-1}{x+1}\right)+2\log_4(x^2-1)\ge2.
\]

\clearpage
\phantomsection\label{task:17}
\section*{Задание № 17}
\subline{Условие}
Электродвижущая сила аккумулятора $E=48$ В, внутреннее сопротивление $r=2$ Ом. Сила тока $I=\dfrac{E}{r+R}$, а мощность на внешнем сопротивлении $P=I^2R$, где $R\ge0$. Найдите наибольшую мощность и соответствующее сопротивление.

\subline{Ваше решение}
Вы правильно составили функцию
\[
P(R)=\frac{48^2R}{(R+2)^2}
\]
и после дифференцирования получили уравнение $-R^2+4=0$. Ниже последовательно записаны $-R^2=-4$, $R=4$ и $R=2$. Последнее значение верное; далее Вы обоснованно получили максимум $288$ Вт и верно исключили отрицательное сопротивление.

\textbf{Ваш итоговый ответ:} $R=2$ Ом, $P_{\max}=288$ Вт.

\subline{Где именно возникает ошибка}
\textbf{Последний правильный шаг:} $-R^2=-4$.

\errorlabel\ Следующая запись $R=4$ математически неверна. После умножения на $-1$ получается $R^2=4$; извлечение квадратного корня даёт два значения $R=\pm2$. Вы затем записали $R=2$, но промежуточный неверный переход остался в решении.

\subline{Почему это неверно}
Если квадрат неизвестного равен четырём, сама неизвестная может равняться как двум, так и минус двум. Число четыре является \emph{значением квадрата}, поэтому непосредственно присваивать его $R$ нельзя. После получения $R=\pm2$ нужно учесть физическое условие $R\ge0$ и оставить $R=2$. Верный конечный ответ не устраняет ошибку в цепочке равенств.

\keyidea\ При переходе от $u^2=a$ к $u=\pm\sqrt a$ контролируйте знак и только затем применяйте область допустимых значений. Для максимума дополнительно исследуйте знак производной.

\subline{Исправление с последнего правильного шага}
\[
-R^2=-4\ \Longrightarrow\ R^2=4\ \Longrightarrow\ R=\pm2.
\]
Так как $R\ge0$, допустимо только $R=2$.

\subline{Полное правильное решение}
\textbf{1. Выразим мощность через сопротивление.}
\[
P(R)=I^2R=\left(\frac{48}{R+2}\right)^2R
=\frac{2304R}{(R+2)^2},\quad R\ge0.
\]
\textbf{2. Найдём производную по правилу производной частного.}
\begin{align*}
P'(R)&=\frac{2304(R+2)^2-2304R\cdot2(R+2)}{(R+2)^4}\\
&=\frac{2304(R+2)[(R+2)-2R]}{(R+2)^4}\\
&=\frac{2304(2-R)}{(R+2)^3}.
\end{align*}
\textbf{3. Исследуем знак.} На $R\ge0$ знаменатель положителен. Поэтому $P'(R)>0$ при $0\le R<2$, $P'(2)=0$ и $P'(R)<0$ при $R>2$. Мощность сначала возрастает, затем убывает, значит, наибольшее значение достигается при $R=2$ Ом.

\textbf{4. Найдём мощность.}
\[
P_{\max}=P(2)=\frac{2304\cdot2}{(2+2)^2}
=\frac{4608}{16}=\boxed{288\text{ Вт}}.
\]
\subline{Проверка}
Для любого $R\ge0$ верно $(R-2)^2\ge0$, откуда $(R+2)^2\ge8R$. При $R>0$ получаем
\[
P(R)=\frac{2304R}{(R+2)^2}\le\frac{2304R}{8R}=288.
\]
Равенство достигается при $R=2$; при $R=0$ мощность равна нулю. Значит, максимум проверен и без производной.

\result\ Максимальная мощность $288$ Вт достигается при $R=2$ Ом. Итоговый ответ был правильным; исправления требует промежуточная алгебраическая запись.

\subline{Задача на закрепление}
Пусть $E=36$ В, $r=3$ Ом, $I=\dfrac{E}{r+R}$ и $P=I^2R$, где $R\ge0$. Найдите наибольшую мощность и сопротивление, при котором она достигается.
\end{document}